\section{大気循環と地衡風・局地風}

\subsection{平成27年度 問24}
\label{q:h27-reg-applied-24}

\noindent\textbf{原典：} sources/H27/17\_applied\_science.pdf、PDF物理9頁。

\textbf{公式正答：} 4（single）。

\paragraph{問題の要旨} 大循環、鉛直層序、大気組成、海陸風、地衡風の記述から適切なものを選ぶ。

\Needspace{4\baselineskip}\paragraph{確認した条件・記述（独自要約）}\mbox{}\par

\noindent\textbf{方向：} 最も適切。\par

\paragraph{問われる核心} 鉛直構造と水平循環を支配する別の条件

\paragraph{必要知識} 鉛直構造と水平循環を支配する別の条件。大循環は緯度、層序は温度勾配、上層組成は混合と拡散、風は力学の条件で判定する。日中の海→陸を述べる4が適切。海陸風を大循環や地衡風の定常近似と同一視しない。

\paragraph{考え方・正答を選ぶ理由（査読待ち初稿）} 大循環は緯度、層序は温度勾配、上層組成は混合と拡散、風は力学の条件で判定する。日中の海→陸を述べる4が適切。海陸風を大循環や地衡風の定常近似と同一視しない。

\Needspace{4\baselineskip}\paragraph{選択肢の検討（査読待ち初稿）}\mbox{}\par

\noindent\textbf{1：不適切} 赤道にモンスーン、極域にハドレー循環。。ハドレーは低緯度の循環であり、極域への割当が違う。モンスーンを低緯度帯全体の基本循環へ一括置換しない。\par

\noindent\textbf{2：不適切} 中間圏オゾンUV加熱で80km温度極大。。オゾン加熱の代表は成層圏で、約50kmに温度極大。中間圏では上方へ冷却し約80〜85kmで極小となる。\par

\noindent\textbf{3：不適切} 成層圏の主成分はO原子とHe。。成層圏でも主成分はN2とO2。O原子・He優勢はさらに高い高度の拡散分離・解離と関係する。\par

\noindent\textbf{4：最適} 日中は地表付近で海から陸へ海風。。日射で陸が海より速く暖まり地表で海→陸の循環を生じる。強い一般風がある場合は局地循環が明瞭でない。\par

\noindent\textbf{5：不適切} 北半球転向力は左向き、地衡風は等圧線直交。。北半球の転向力は進行方向の右。地衡風は気圧傾度力と釣り合い等圧線にほぼ沿う。\par

\paragraph{再利用する考え方} 全球平均・地域差、貯留量・流量、自然温室効果・追加強制力、代表値・普遍則を区別する。鉛直構造と水平循環を支配する別の条件の対象と成立条件を確認する。

\paragraph{共通理論} \hyperref[theory:atmosphere:layers]{第\ref*{theory:atmosphere:layers}節}, \hyperref[theory:atmosphere:wind]{第\ref*{theory:atmosphere:wind}節}, \hyperref[theory:atmosphere:local]{第\ref*{theory:atmosphere:local}節}

\par 参考：\cite{atm_noaa_layers,atm_noaa_winds,atm_jma_dynamic}。

\paragraph{関連問題} \hyperref[q:r02-reg-applied-21]{R02-REG-APPLIED-21}

\par 難易度 2（校準初稿）、優先度未評価。

\subsection{令和2年度 問21}
\label{q:r02-reg-applied-21}

\noindent\textbf{原典：} sources/R02/17\_applied\_science.pdf、PDF物理8頁。

\textbf{公式正答：} 4（single）。

\par\textbf{検討未確定：} 肢3の170km超O・1000km付近He主要の数値高度は一次数表未確認。原典の数字は保持し部分校正。

\paragraph{問題の要旨} 大循環、成層圏加熱、上層組成、海陸風と地衡風の記述を比較する。

\Needspace{4\baselineskip}\paragraph{確認した条件・記述（独自要約）}\mbox{}\par

\noindent\textbf{方向：} 最も不適切。\par

\paragraph{問われる核心} 鉛直組成と循環モデルの適用条件

\paragraph{必要知識} 鉛直組成と循環モデルの適用条件。日中の陸加熱と夜間の陸冷却を分けると夜間海→陸の4が不適切となる。大循環は緯度帯、地衡風は運動方程式の近似で判定する。3の原子O・Heの数値高度は定性的な分離の理解とは別に一次数表を要するため、全肢確認へ昇格させない。

\paragraph{考え方・正答を選ぶ理由（査読待ち初稿）} 日中の陸加熱と夜間の陸冷却を分けると夜間海→陸の4が不適切となる。大循環は緯度帯、地衡風は運動方程式の近似で判定する。3の原子O・Heの数値高度は定性的な分離の理解とは別に一次数表を要するため、全肢確認へ昇格させない。

\Needspace{4\baselineskip}\paragraph{選択肢の検討（査読待ち初稿）}\mbox{}\par

\noindent\textbf{1：適切（平均模式）} 低緯度ハドレー、中緯度フェレル、高緯度極循環。。緯度帯に対応する帯状平均循環の基本模式である。瞬時の風を3つの閉じた固定セルだけで表さない。\par

\noindent\textbf{2：適切} 成層圏オゾンUV吸収で約50km温度極大。。層序とオゾンの放射作用が整合する。\par

\noindent\textbf{3：数値高度は要確認} 80kmまでほぼ一定、170km超でO原子、1000km付近でHe主要。。均質圏から上方の解離と拡散分離の定性的説明は合う。優勢成分が切り替わる高度は太陽活動等に依存し、この2数値の一次数表検証は未完了。\par

\noindent\textbf{4：不適切} 夜間は地表で海から陸へ海風。。夜間は陸が速く冷え、基本の地表循環は陸→海の陸風。\par

\noindent\textbf{5：適切（近似条件）} 傾度力と転向力の釣合いで等圧線沿いの地衡風。。摩擦・加速度が小さく赤道から離れた条件での理論風である。\par

\paragraph{再利用する考え方} 全球平均・地域差、貯留量・流量、自然温室効果・追加強制力、代表値・普遍則を区別する。鉛直組成と循環モデルの適用条件の対象と成立条件を確認する。

\paragraph{共通理論} \hyperref[theory:atmosphere:composition]{第\ref*{theory:atmosphere:composition}節}, \hyperref[theory:atmosphere:wind]{第\ref*{theory:atmosphere:wind}節}, \hyperref[theory:atmosphere:local]{第\ref*{theory:atmosphere:local}節}

\par 参考：\cite{atm_noaa_winds,atm_noaa_layers,atm_us1976}。

\paragraph{関連問題} \hyperref[q:h27-reg-applied-24]{H27-REG-APPLIED-24}

\par 難易度 2（校準初稿）、優先度未評価。

\section{乾燥大気の組成と微量気体の循環}

\subsection{平成24年度 問34}
\label{q:h24-reg-applied-34}

\noindent\textbf{原典：} sources/H24/17\_applied\_science.pdf、PDF物理10頁。

\textbf{公式正答：} 2（single）。

\paragraph{問題の要旨} 大気微量成分の発生源、温室効果、成層圏オゾンへの作用を比較する。

\Needspace{4\baselineskip}\paragraph{確認した条件・記述（独自要約）}\mbox{}\par

\noindent\textbf{方向：} 最も不適切。\par

\paragraph{問われる核心} 微量気体の源・寿命・放射と化学反応

\paragraph{必要知識} 微量気体の源・寿命・放射と化学反応。大気成分の量と放射・化学上の作用は別である。微量でも温室効果やオゾン反応へ寄与する。メタンの源は化石燃料採掘のほか湿地、水田、家畜等に広がるため2が不適切。N2OをN2やNO2と取り違えず、オゾンのUV作用と温室気体の赤外作用を区別する。

\paragraph{考え方・正答を選ぶ理由（査読待ち初稿）} 大気成分の量と放射・化学上の作用は別である。微量でも温室効果やオゾン反応へ寄与する。メタンの源は化石燃料採掘のほか湿地、水田、家畜等に広がるため2が不適切。N2OをN2やNO2と取り違えず、オゾンのUV作用と温室気体の赤外作用を区別する。

\Needspace{4\baselineskip}\paragraph{選択肢の検討（査読待ち初稿）}\mbox{}\par

\noindent\textbf{1：適切（条件付き概数）} N2Oは温室効果とオゾン破壊に関与。。N2Oから生成されるNO等がオゾンの触媒反応へ関わり、赤外吸収も持つ。\par

\noindent\textbf{2：不適切} メタンは湿地・水田・反すう動物の寄与が小さく採掘由来がほぼ全て。。嫌気的な微生物過程や家畜・水田なども重要な源で、化石燃料採掘だけに限定できない。\par

\noindent\textbf{3：適切（条件付き概数）} CO2は季節変動を伴い長期増加。。原典当時の観測と整合。季節周期の上へ長期増加が重なる。\par

\noindent\textbf{4：適切（条件付き概数）} CFC11/12は安定で輸送後、成層圏でUV分解しClを放出。。対流圏で壊れにくい性質が成層圏への輸送を可能にし、UVで遊離した塩素が反応へ関与する。\par

\noindent\textbf{5：適切（条件付き概数）} 成層圏オゾンは有害UVを吸収し大気を加熱。。オゾンの吸収が成層圏の正の温度勾配を作る要因である。\par

\paragraph{再利用する考え方} 全球平均・地域差、貯留量・流量、自然温室効果・追加強制力、代表値・普遍則を区別する。微量気体の源・寿命・放射と化学反応の対象と成立条件を確認する。

\paragraph{共通理論} \hyperref[theory:atmosphere:composition]{第\ref*{theory:atmosphere:composition}節}, \hyperref[theory:atmosphere:ozone]{第\ref*{theory:atmosphere:ozone}節}

\par 参考：\cite{atm_noaa_ghg,atm_jma_ozone}。

\par 難易度 2（校準初稿）、優先度未評価。

\subsection{平成30年度 問20}
\label{q:h30-reg-applied-20}

\noindent\textbf{原典：} sources/H30/17\_applied\_science.pdf、PDF物理7頁。

\textbf{公式正答：} 4（single）。

\paragraph{問題の要旨} 乾燥大気の主要成分の量順と、微量温室気体の発生・季節変動を判定する。

\Needspace{4\baselineskip}\paragraph{確認した条件・記述（独自要約）}\mbox{}\par

\noindent\textbf{方向：} 最も不適切。\par

\paragraph{問われる核心} 乾燥大気の混合比と変動性の違い

\paragraph{必要知識} 乾燥大気の混合比と変動性の違い。乾燥大気のN2約78\%、O2約21\%、Ar約0.93\%の順を先に確かめる。CO2はArより少ないので順位を入れ替えた肢が不適切。微量でも放射・生物循環の作用は大きく、体積比の多寡から温室作用の大きさを直接決めない。

\paragraph{考え方・正答を選ぶ理由（査読待ち初稿）} 乾燥大気のN2約78\%、O2約21\%、Ar約0.93\%の順を先に確かめる。CO2はArより少ないので順位を入れ替えた肢が不適切。微量でも放射・生物循環の作用は大きく、体積比の多寡から温室作用の大きさを直接決めない。

\Needspace{4\baselineskip}\paragraph{選択肢の検討（査読待ち初稿）}\mbox{}\par

\noindent\textbf{1：適切（条件付き概数）} N2は安定で全球の混合比がほぼ一様。。よく混合された下層大気の乾燥混合比の代表。局所的な反応が全くないという意味ではない。\par

\noindent\textbf{2：適切（条件付き概数）} O2は原始大気に乏しく地史で増加。。酸素発生型光合成等と地球表層の酸化還元過程の長期変化に関係する。現在の組成を原始大気へ遡及しない。\par

\noindent\textbf{3：適切（条件付き概数）} CO2は光合成・呼吸で季節変動。。特に北半球の植生季節活動が観測曲線へ現れる。長期人為増加とは別成分。\par

\noindent\textbf{4：不適切} ArはN2・O2・CO2の次に多い。。乾燥大気の量順はN2、O2、Ar、CO2。Ar約0.93\%はCO2の約0.04\%より大きい。\par

\noindent\textbf{5：適切（条件付き概数）} CH4は嫌気的条件で生成する温室気体で近年増加。。生物起源の生成は酸素の乏しい環境で起こる。源は多様で、ここでいう増加は問題出題の近年に関する。\par

\paragraph{再利用する考え方} 全球平均・地域差、貯留量・流量、自然温室効果・追加強制力、代表値・普遍則を区別する。乾燥大気の混合比と変動性の違いの対象と成立条件を確認する。

\paragraph{共通理論} \hyperref[theory:atmosphere:composition]{第\ref*{theory:atmosphere:composition}節}

\par 参考：\cite{atm_noaa_composition,atm_noaa_ghg}。

\paragraph{関連問題} \hyperref[q:r05-reg-applied-18]{R05-REG-APPLIED-18}

\par 難易度 2（校準初稿）、優先度未評価。

\subsection{令和5年度 問18}
\label{q:r05-reg-applied-18}

\noindent\textbf{原典：} sources/R05/17\_applied\_science.pdf、PDF物理6頁。

\textbf{公式正答：} 1（single）。

\paragraph{問題の要旨} 乾燥大気の主要成分の量順と、微量温室気体の発生・季節変動を判定する。

\Needspace{4\baselineskip}\paragraph{確認した条件・記述（独自要約）}\mbox{}\par

\noindent\textbf{方向：} 最も不適切。\par

\paragraph{問われる核心} 乾燥大気の混合比と変動性の違い

\paragraph{必要知識} 乾燥大気の混合比と変動性の違い。乾燥大気のN2約78\%、O2約21\%、Ar約0.93\%の順を先に確かめる。CO2はArより少ないので順位を入れ替えた肢が不適切。微量でも放射・生物循環の作用は大きく、体積比の多寡から温室作用の大きさを直接決めない。

\paragraph{考え方・正答を選ぶ理由（査読待ち初稿）} 乾燥大気のN2約78\%、O2約21\%、Ar約0.93\%の順を先に確かめる。CO2はArより少ないので順位を入れ替えた肢が不適切。微量でも放射・生物循環の作用は大きく、体積比の多寡から温室作用の大きさを直接決めない。

\Needspace{4\baselineskip}\paragraph{選択肢の検討（査読待ち初稿）}\mbox{}\par

\noindent\textbf{1：不適切} ArはN2・O2・CO2の次に多い。。乾燥大気の量順はN2、O2、Ar、CO2。Ar約0.93\%はCO2の約0.04\%より大きい。\par

\noindent\textbf{2：適切（条件付き概数）} CH4は嫌気的条件で生成する温室気体で近年増加。。生物起源の生成は酸素の乏しい環境で起こる。源は多様で、ここでいう増加は問題出題の近年に関する。\par

\noindent\textbf{3：適切（条件付き概数）} N2は安定で全球の混合比がほぼ一様。。よく混合された下層大気の乾燥混合比の代表。局所的な反応が全くないという意味ではない。\par

\noindent\textbf{4：適切（条件付き概数）} O2は原始大気に乏しく地史で増加。。酸素発生型光合成等と地球表層の酸化還元過程の長期変化に関係する。現在の組成を原始大気へ遡及しない。\par

\noindent\textbf{5：適切（条件付き概数）} CO2は光合成・呼吸で季節変動。。特に北半球の植生季節活動が観測曲線へ現れる。長期人為増加とは別成分。\par

\paragraph{再利用する考え方} 全球平均・地域差、貯留量・流量、自然温室効果・追加強制力、代表値・普遍則を区別する。乾燥大気の混合比と変動性の違いの対象と成立条件を確認する。

\paragraph{共通理論} \hyperref[theory:atmosphere:composition]{第\ref*{theory:atmosphere:composition}節}

\par 参考：\cite{atm_noaa_composition,atm_noaa_ghg}。

\paragraph{関連問題} \hyperref[q:h30-reg-applied-20]{H30-REG-APPLIED-20}

\par 難易度 2（校準初稿）、優先度未評価。

\section{大気の鉛直温度・密度構造}

\subsection{令和元年度 問19}
\label{q:r01-reg-applied-19}

\noindent\textbf{原典：} sources/R01/17\_applied\_science.pdf、PDF物理6頁。

\textbf{公式正答：} 4（single）。

\par\textbf{検討未確定：} 肢3の10km密度1/4と標準モデル約0.34に差。原典密度を圧力へ改変しない。概数の許容精度をAstra査読待ち。

\paragraph{問題の要旨} 大気の層序、組成、密度低下、成層圏加熱と均質圏を比較する。

\Needspace{4\baselineskip}\paragraph{確認した条件・記述（独自要約）}\mbox{}\par

\noindent\textbf{方向：} 最も不適切。\par

\paragraph{問われる核心} 鉛直温度構造と質量密度・混合比を区別する

\paragraph{必要知識} 鉛直温度構造と質量密度・混合比を区別する。50km付近の温度極大をCO2の赤外吸収に帰す4が不適切である。オゾンのUV吸収で成層圏を加熱する。ただし3の密度1/4は標準モデル約0.34との精度差を残し、圧力への無断校正をしない。公式4と独立4の一致は全ての概数の厳密保証ではない。

\paragraph{考え方・正答を選ぶ理由（査読待ち初稿）} 50km付近の温度極大をCO2の赤外吸収に帰す4が不適切である。オゾンのUV吸収で成層圏を加熱する。ただし3の密度1/4は標準モデル約0.34との精度差を残し、圧力への無断校正をしない。公式4と独立4の一致は全ての概数の厳密保証ではない。

\Needspace{4\baselineskip}\paragraph{選択肢の検討（査読待ち初稿）}\mbox{}\par

\noindent\textbf{1：適切} 下から対流圏・成層圏・中間圏・熱圏。。主として気温の鉛直勾配を基準とする区分である。\par

\noindent\textbf{2：適切（概数）} 空気は体積でN2約8割、O2約2割。。端数の微量成分や水蒸気を除いた粗い概数として適切。\par

\noindent\textbf{3：概数精度に疑義} 10kmで質量密度は地表の約1/4。。原典は密度と記載。標準モデルで約0.414kg/m\textsuperscript{3}対地表1.225、比約0.338となり1/4とは差がある。圧力比約0.26との混同があり得るが、原典を圧力へ改変しない。\par

\noindent\textbf{4：不適切} 50km付近の温度極大はCO2赤外吸収による。。代表的な成層圏界面の加熱はオゾンの紫外吸収。温室気体一般の赤外作用へ置換できない。\par

\noindent\textbf{5：適切（主要乾燥成分）} 乾燥主要成分比は80km程度までほぼ一定。。乱流混合が分子拡散より強い均質圏の主要長寿命気体を指す。水蒸気・オゾン等まで一様とはしない。\par

\paragraph{再利用する考え方} 全球平均・地域差、貯留量・流量、自然温室効果・追加強制力、代表値・普遍則を区別する。鉛直温度構造と質量密度・混合比を区別するの対象と成立条件を確認する。

\paragraph{共通理論} \hyperref[theory:atmosphere:layers]{第\ref*{theory:atmosphere:layers}節}, \hyperref[theory:atmosphere:hydrostatic]{第\ref*{theory:atmosphere:hydrostatic}節}, \hyperref[theory:atmosphere:ozone]{第\ref*{theory:atmosphere:ozone}節}

\par 参考：\cite{atm_noaa_layers,atm_us1976,atm_jma_ozone}。

\paragraph{関連問題} \hyperref[q:r07-reg-applied-21]{R07-REG-APPLIED-21}

\par 難易度 2（校準初稿）、優先度未評価。

\subsection{令和元年度再試験 問21}
\label{q:r01-ret-applied-21}

\noindent\textbf{原典：} sources/R01\_retest/17\_applied\_science.pdf、PDF物理7頁。

\textbf{公式正答：} 2（single）。

\paragraph{問題の要旨} 対流圏の質量・気温減率・界面高度と成層圏温度、オゾンの分布を判定する。

\Needspace{4\baselineskip}\paragraph{確認した条件・記述（独自要約）}\mbox{}\par

\noindent\textbf{方向：} 最も不適切。\par

\paragraph{問われる核心} 大気の鉛直構造と環境・断熱気温減率

\paragraph{必要知識} 大気の鉛直構造と環境・断熱気温減率。環境減率の標準値は約6.5K/kmであり、3K/kmという肢が不適切。温度の観測平均と、気塊の断熱移動を表す9.8K/km等を混ぜない。界面は赤道側で高く、オゾンが多い成層圏ではUV加熱による温度上昇が生じる。

\paragraph{考え方・正答を選ぶ理由（査読待ち初稿）} 環境減率の標準値は約6.5K/kmであり、3K/kmという肢が不適切。温度の観測平均と、気塊の断熱移動を表す9.8K/km等を混ぜない。界面は赤道側で高く、オゾンが多い成層圏ではUV加熱による温度上昇が生じる。

\Needspace{4\baselineskip}\paragraph{選択肢の検討（査読待ち初稿）}\mbox{}\par

\noindent\textbf{1：適切（条件付き概数）} 対流圏質量は気圏の約3/4。。大気質量が下層へ集中する代表的な全球概数。界面の緯度・季節差で厳密値は変わる。\par

\noindent\textbf{2：不適切} 対流圏の気温減率は3℃/km。。標準的な対流圏環境減率は約6.5K/km。乾燥断熱9.8K/kmや湿潤断熱・実際の逆転層と区別する。\par

\noindent\textbf{3：適切（条件付き概数）} 対流圏界面は低緯度で高く高緯度で低い。。低緯度の高い温度と深い対流により界面は高くなる。\par

\noindent\textbf{4：適切（条件付き概数）} 成層圏では高度とともに気温上昇。。オゾン吸収による加熱が上方の温度増加を作る。成層圏の下部がほぼ等温の場合を無視した厳密な全高度単調則ではない。\par

\noindent\textbf{5：適切（条件付き概数）} オゾンの約9割は成層圏。。全球オゾン鉛直分布の代表値。地表オゾン汚染の存在を否定しない。\par

\paragraph{再利用する考え方} 全球平均・地域差、貯留量・流量、自然温室効果・追加強制力、代表値・普遍則を区別する。大気の鉛直構造と環境・断熱気温減率の対象と成立条件を確認する。

\paragraph{共通理論} \hyperref[theory:atmosphere:layers]{第\ref*{theory:atmosphere:layers}節}, \hyperref[theory:atmosphere:lapse]{第\ref*{theory:atmosphere:lapse}節}, \hyperref[theory:atmosphere:ozone]{第\ref*{theory:atmosphere:ozone}節}

\par 参考：\cite{atm_noaa_layers,atm_us1976,atm_jma_ozone}。

\paragraph{関連問題} \hyperref[q:r04-reg-applied-22]{R04-REG-APPLIED-22}

\par 難易度 2（校準初稿）、優先度未評価。

\subsection{令和4年度 問22}
\label{q:r04-reg-applied-22}

\noindent\textbf{原典：} sources/R04/17\_applied\_science.pdf、PDF物理7頁。

\textbf{公式正答：} 1（single）。

\paragraph{問題の要旨} 対流圏の質量・気温減率・界面高度と成層圏温度、オゾンの分布を判定する。

\Needspace{4\baselineskip}\paragraph{確認した条件・記述（独自要約）}\mbox{}\par

\noindent\textbf{方向：} 最も不適切。\par

\paragraph{問われる核心} 大気の鉛直構造と環境・断熱気温減率

\paragraph{必要知識} 大気の鉛直構造と環境・断熱気温減率。環境減率の標準値は約6.5K/kmであり、3K/kmという肢が不適切。温度の観測平均と、気塊の断熱移動を表す9.8K/km等を混ぜない。界面は赤道側で高く、オゾンが多い成層圏ではUV加熱による温度上昇が生じる。

\paragraph{考え方・正答を選ぶ理由（査読待ち初稿）} 環境減率の標準値は約6.5K/kmであり、3K/kmという肢が不適切。温度の観測平均と、気塊の断熱移動を表す9.8K/km等を混ぜない。界面は赤道側で高く、オゾンが多い成層圏ではUV加熱による温度上昇が生じる。

\Needspace{4\baselineskip}\paragraph{選択肢の検討（査読待ち初稿）}\mbox{}\par

\noindent\textbf{1：不適切} 対流圏の気温減率は3℃/km。。標準的な対流圏環境減率は約6.5K/km。乾燥断熱9.8K/kmや湿潤断熱・実際の逆転層と区別する。\par

\noindent\textbf{2：適切（条件付き概数）} オゾンの約9割は成層圏。。全球オゾン鉛直分布の代表値。地表オゾン汚染の存在を否定しない。\par

\noindent\textbf{3：適切（条件付き概数）} 成層圏では高度とともに気温上昇。。オゾン吸収による加熱が上方の温度増加を作る。成層圏の下部がほぼ等温の場合を無視した厳密な全高度単調則ではない。\par

\noindent\textbf{4：適切（条件付き概数）} 対流圏界面は低緯度で高く高緯度で低い。。低緯度の高い温度と深い対流により界面は高くなる。\par

\noindent\textbf{5：適切（条件付き概数）} 対流圏質量は気圏の約3/4。。大気質量が下層へ集中する代表的な全球概数。界面の緯度・季節差で厳密値は変わる。\par

\paragraph{再利用する考え方} 全球平均・地域差、貯留量・流量、自然温室効果・追加強制力、代表値・普遍則を区別する。大気の鉛直構造と環境・断熱気温減率の対象と成立条件を確認する。

\paragraph{共通理論} \hyperref[theory:atmosphere:layers]{第\ref*{theory:atmosphere:layers}節}, \hyperref[theory:atmosphere:lapse]{第\ref*{theory:atmosphere:lapse}節}, \hyperref[theory:atmosphere:ozone]{第\ref*{theory:atmosphere:ozone}節}

\par 参考：\cite{atm_noaa_layers,atm_us1976,atm_jma_ozone}。

\paragraph{関連問題} \hyperref[q:r01-ret-applied-21]{R01-RET-APPLIED-21}

\par 難易度 2（校準初稿）、優先度未評価。

\subsection{令和7年度 問21}
\label{q:r07-reg-applied-21}

\noindent\textbf{原典：} sources/R07/17\_applied\_science.pdf、PDF物理7頁。

\textbf{公式正答：} 2（single）。

\par\textbf{検討未確定：} 肢1の10km密度1/4は標準モデル約0.34と差があり概数精度を査読待ち。原典の数値・密度という語は保持。

\paragraph{問題の要旨} 大気密度の概数、界面の緯度差、主要成分の混合、オゾン分布と層序を比較する。

\Needspace{4\baselineskip}\paragraph{確認した条件・記述（独自要約）}\mbox{}\par

\noindent\textbf{方向：} 最も不適切。\par

\paragraph{問われる核心} 温度層序・密度と緯度変化

\paragraph{必要知識} 温度層序・密度と緯度変化。低緯度の対流圏界面は高く、極域は低いので2が不適切。ただし1は標準モデルの10km密度比約0.34と1/4の精度差があり、原典の密度を圧力へ書き換えず部分校正に残す。公式正答2と別肢の概数精度の査読を分ける。

\paragraph{考え方・正答を選ぶ理由（査読待ち初稿）} 低緯度の対流圏界面は高く、極域は低いので2が不適切。ただし1は標準モデルの10km密度比約0.34と1/4の精度差があり、原典の密度を圧力へ書き換えず部分校正に残す。公式正答2と別肢の概数精度の査読を分ける。

\Needspace{4\baselineskip}\paragraph{選択肢の検討（査読待ち初稿）}\mbox{}\par

\noindent\textbf{1：概数精度に疑義} 10kmで質量密度は地表の約1/4。。原典は密度と記載。標準モデルで約0.414kg/m\textsuperscript{3}対地表1.225、比約0.338となり1/4とは差がある。圧力比約0.26との混同があり得るが、原典を圧力へ改変しない。\par

\noindent\textbf{2：不適切} 対流圏界面は赤道より極域で高い。。低緯度では深い対流と高温により界面は高い。逆にしたこの肢が明白な誤り。\par

\noindent\textbf{3：適切（主要成分）} N2/O2/Arの組成は数10kmまでほぼ一定。。均質圏の主要長寿命成分の混合比の概数。微量気体や水蒸気の分布は同じではない。\par

\noindent\textbf{4：適切（概数）} 全オゾンの約10\%が対流圏。。約90\%成層圏という代表分布の補集合であり整合する。\par

\noindent\textbf{5：適切} 下から対流圏・成層圏・中間圏・熱圏。。主として気温の鉛直勾配を基準とする区分である。\par

\paragraph{再利用する考え方} 全球平均・地域差、貯留量・流量、自然温室効果・追加強制力、代表値・普遍則を区別する。温度層序・密度と緯度変化の対象と成立条件を確認する。

\paragraph{共通理論} \hyperref[theory:atmosphere:layers]{第\ref*{theory:atmosphere:layers}節}, \hyperref[theory:atmosphere:hydrostatic]{第\ref*{theory:atmosphere:hydrostatic}節}, \hyperref[theory:atmosphere:ozone]{第\ref*{theory:atmosphere:ozone}節}

\par 参考：\cite{atm_noaa_layers,atm_us1976,atm_jma_ozone}。

\paragraph{関連問題} \hyperref[q:r01-reg-applied-19]{R01-REG-APPLIED-19}

\par 難易度 2（校準初稿）、優先度未評価。

\section{地球の受光幾何と放射収支}

\subsection{平成26年度 問26}
\label{q:h26-reg-applied-26}

\noindent\textbf{原典：} sources/H26/17\_applied\_science.pdf、PDF物理11頁。

\textbf{公式正答：} 1（single）。

\paragraph{問題の要旨} 地球の反射率、太陽放射帯域、季節原因、太陽定数と緯度間の熱収支を判定する。

\Needspace{4\baselineskip}\paragraph{確認した条件・記述（独自要約）}\mbox{}\par

\noindent\textbf{方向：} 最も適切。\par

\paragraph{問われる核心} 短波入射と長波放射の地球規模の収支

\paragraph{必要知識} 短波入射と長波放射の地球規模の収支。反射と吸収、短波日射と長波地球放射を分ける。反射の約31\%が問題採用の適切な概数である。四季は地軸傾斜、太陽定数は大気上端の垂直面、放射不均衡は低緯度から高緯度への輸送につながる。31\%を現在の精密観測で必ず同じとするのではなく、問題の比較精度で扱う。

\paragraph{考え方・正答を選ぶ理由（査読待ち初稿）} 反射と吸収、短波日射と長波地球放射を分ける。反射の約31\%が問題採用の適切な概数である。四季は地軸傾斜、太陽定数は大気上端の垂直面、放射不均衡は低緯度から高緯度への輸送につながる。31\%を現在の精密観測で必ず同じとするのではなく、問題の比較精度で扱う。

\Needspace{4\baselineskip}\paragraph{選択肢の検討（査読待ち初稿）}\mbox{}\par

\noindent\textbf{1：最適} 入射の約31\%が雲・大気・地表で反射散乱され宇宙へ戻る。。問題で採る全球平均反射率約3割の代表値と合う。時間・雲・地表で変化するので正確な不変定数ではない。\par

\noindent\textbf{2：不適切} 太陽放射の約半分はガンマ線・X線。。太陽放射の主な帯域は可視・近赤外であり、高エネルギー放射が半分ではない。\par

\noindent\textbf{3：不適切} 地球太陽距離の変化が四季の原因。。地軸の傾きで太陽高度と日長が変わることが主因であり、南北半球の季節は逆になる。\par

\noindent\textbf{4：不適切} 大気下端の太陽垂直面の受光を太陽定数という。。大気の吸収散乱の影響を除いた大気上端で、太陽光に垂直な単位面積あたりの放射束で定義する。\par

\noindent\textbf{5：不適切} 吸収日射の緯度差は小さく、高緯度の放出が低緯度より大きい。。低緯度の日射吸収は大きく、正味収支の緯度差を大気・海洋の熱輸送が補う。\par

\paragraph{再利用する考え方} 全球平均・地域差、貯留量・流量、自然温室効果・追加強制力、代表値・普遍則を区別する。短波入射と長波放射の地球規模の収支の対象と成立条件を確認する。

\paragraph{共通理論} \hyperref[theory:atmosphere:radiation]{第\ref*{theory:atmosphere:radiation}節}, \hyperref[theory:atmosphere:transport]{第\ref*{theory:atmosphere:transport}節}

\par 参考：\cite{atm_nasa_budget,atm_nasa_albedo,atm_nasa_transport}。

\paragraph{関連問題} \hyperref[q:r01-reg-applied-22]{R01-REG-APPLIED-22}, \hyperref[q:r06-reg-applied-20]{R06-REG-APPLIED-20}

\par 難易度 2（校準初稿）、優先度未評価。

\subsection{平成28年度 問26}
\label{q:h28-reg-applied-26}

\noindent\textbf{原典：} sources/H28/17\_applied\_science.pdf、PDF物理11頁。

\textbf{公式正答：} 1（single）。

\paragraph{問題の要旨} 指定された地球半径と太陽定数から、地球が毎秒受け取る日射の総量を求める。

\Needspace{4\baselineskip}\paragraph{確認した条件・記述（独自要約）}\mbox{}\par

\noindent\textbf{方向：} 最も近い値。\par

\noindent\textbf{半径：} 6370km=\allowbreak{}6.370×10\textasciicircum{}6m。\par

\noindent\textbf{太陽定数：} 1.368×10\textasciicircum{}3W/m\textsuperscript{2}。\par

\noindent\textbf{大気：} H28は反射吸収無視、H30は大気上端で定義。\par

\paragraph{問われる核心} 地球の受光断面と球面平均を区別する

\paragraph{必要知識} 地球の受光断面と球面平均を区別する。平行入射光を遮る投影面積は\(\pi R^2\)で、球表面積ではない。\(P=S_0\pi R^2=1368\pi(6.370\times10^6)^2\simeq1.744\times10^{17}\ {\rm W}\)。Wは毎秒のエネルギーを表す。全球平均面積当たりならこの総量を\(4\pi R^2\)で割り\(S_0/4\)となるが、今回は総量を問う。

\paragraph{考え方・正答を選ぶ理由（査読待ち初稿）} 平行入射光を遮る投影面積は\(\pi R^2\)で、球表面積ではない。\(P=S_0\pi R^2=1368\pi(6.370\times10^6)^2\simeq1.744\times10^{17}\ {\rm W}\)。Wは毎秒のエネルギーを表す。全球平均面積当たりならこの総量を\(4\pi R^2\)で割り\(S_0/4\)となるが、今回は総量を問う。

\Needspace{4\baselineskip}\paragraph{選択肢の検討（査読待ち初稿）}\mbox{}\par

\noindent\textbf{1：最適} \(1.74\times10^{17}\ {\rm W}\)。受光円の計算1.74×10\textasciicircum{}17Wと一致。\par

\noindent\textbf{2：不適切} \(3.49\times10^{17}\ {\rm W}\)。計算値の約2倍。受光断面と球表面を混同せず、単位換算と面積を確認する。\par

\noindent\textbf{3：不適切} \(6.97\times10^{17}\ {\rm W}\)。計算値の約4倍。受光断面と球表面を混同せず、単位換算と面積を確認する。\par

\noindent\textbf{4：不適切} \(1.39\times10^{18}\ {\rm W}\)。計算値の約8倍。受光断面と球表面を混同せず、単位換算と面積を確認する。\par

\noindent\textbf{5：不適切} \(2.79\times10^{18}\ {\rm W}\)。計算値の約16倍。受光断面と球表面を混同せず、単位換算と面積を確認する。\par

\paragraph{再利用する考え方} 全球平均・地域差、貯留量・流量、自然温室効果・追加強制力、代表値・普遍則を区別する。地球の受光断面と球面平均を区別するの対象と成立条件を確認する。

\paragraph{共通理論} \hyperref[theory:atmosphere:geometry]{第\ref*{theory:atmosphere:geometry}節}

\par 参考：\cite{atm_nasa_budget}。

\paragraph{関連問題} \hyperref[q:h30-reg-applied-22]{H30-REG-APPLIED-22}, \hyperref[q:r01-ret-applied-20]{R01-RET-APPLIED-20}

\par 難易度 2（校準初稿）、優先度未評価。

\subsection{平成30年度 問22}
\label{q:h30-reg-applied-22}

\noindent\textbf{原典：} sources/H30/17\_applied\_science.pdf、PDF物理7頁。

\textbf{公式正答：} 2（single）。

\paragraph{問題の要旨} 指定された地球半径と太陽定数から、地球が毎秒受け取る日射の総量を求める。

\Needspace{4\baselineskip}\paragraph{確認した条件・記述（独自要約）}\mbox{}\par

\noindent\textbf{方向：} 最も近い値。\par

\noindent\textbf{半径：} 6370km=\allowbreak{}6.370×10\textasciicircum{}6m。\par

\noindent\textbf{太陽定数：} 1.368×10\textasciicircum{}3W/m\textsuperscript{2}。\par

\noindent\textbf{大気：} H28は反射吸収無視、H30は大気上端で定義。\par

\paragraph{問われる核心} 地球の受光断面と球面平均を区別する

\paragraph{必要知識} 地球の受光断面と球面平均を区別する。平行入射光を遮る投影面積は\(\pi R^2\)で、球表面積ではない。\(P=S_0\pi R^2=1368\pi(6.370\times10^6)^2\simeq1.744\times10^{17}\ {\rm W}\)。Wは毎秒のエネルギーを表す。全球平均面積当たりならこの総量を\(4\pi R^2\)で割り\(S_0/4\)となるが、今回は総量を問う。

\paragraph{考え方・正答を選ぶ理由（査読待ち初稿）} 平行入射光を遮る投影面積は\(\pi R^2\)で、球表面積ではない。\(P=S_0\pi R^2=1368\pi(6.370\times10^6)^2\simeq1.744\times10^{17}\ {\rm W}\)。Wは毎秒のエネルギーを表す。全球平均面積当たりならこの総量を\(4\pi R^2\)で割り\(S_0/4\)となるが、今回は総量を問う。

\Needspace{4\baselineskip}\paragraph{選択肢の検討（査読待ち初稿）}\mbox{}\par

\noindent\textbf{1：不適切} \(8.71\times10^{16}\ {\rm W}\)。計算値の約0.5倍。受光断面と球表面を混同せず、単位換算と面積を確認する。\par

\noindent\textbf{2：最適} \(1.74\times10^{17}\ {\rm W}\)。受光円の計算1.74×10\textasciicircum{}17Wと一致。\par

\noindent\textbf{3：不適切} \(3.49\times10^{17}\ {\rm W}\)。計算値の約2倍。受光断面と球表面を混同せず、単位換算と面積を確認する。\par

\noindent\textbf{4：不適切} \(6.97\times10^{17}\ {\rm W}\)。計算値の約4倍。受光断面と球表面を混同せず、単位換算と面積を確認する。\par

\noindent\textbf{5：不適切} \(1.39\times10^{18}\ {\rm W}\)。計算値の約8倍。受光断面と球表面を混同せず、単位換算と面積を確認する。\par

\paragraph{再利用する考え方} 全球平均・地域差、貯留量・流量、自然温室効果・追加強制力、代表値・普遍則を区別する。地球の受光断面と球面平均を区別するの対象と成立条件を確認する。

\paragraph{共通理論} \hyperref[theory:atmosphere:geometry]{第\ref*{theory:atmosphere:geometry}節}

\par 参考：\cite{atm_nasa_budget}。

\paragraph{関連問題} \hyperref[q:h28-reg-applied-26]{H28-REG-APPLIED-26}

\par 難易度 2（校準初稿）、優先度未評価。

\subsection{令和元年度 問22}
\label{q:r01-reg-applied-22}

\noindent\textbf{原典：} sources/R01/17\_applied\_science.pdf、PDF物理7頁。

\textbf{公式正答：} 3（single）。

\paragraph{問題の要旨} 季節・緯度収支、太陽放射の主帯域、反射率、太陽定数の定義を比較する。

\Needspace{4\baselineskip}\paragraph{確認した条件・記述（独自要約）}\mbox{}\par

\noindent\textbf{方向：} 最も適切。\par

\paragraph{問われる核心} 放射の波長帯・反射率・幾何の区別

\paragraph{必要知識} 放射の波長帯・反射率・幾何の区別。太陽放射の波長分布は可視・近赤外が中心である。3の約半分可視が適切。他肢の距離を季節原因とする説明、緯度収支の逆転、15\%反射、大気下端定義を分けて否定する。

\paragraph{考え方・正答を選ぶ理由（査読待ち初稿）} 太陽放射の波長分布は可視・近赤外が中心である。3の約半分可視が適切。他肢の距離を季節原因とする説明、緯度収支の逆転、15\%反射、大気下端定義を分けて否定する。

\Needspace{4\baselineskip}\paragraph{選択肢の検討（査読待ち初稿）}\mbox{}\par

\noindent\textbf{1：不適切} 地球太陽距離の変化が四季の原因。。地軸の傾きで太陽高度と日長が変わることが主因であり、南北半球の季節は逆になる。\par

\noindent\textbf{2：不適切} 吸収日射の緯度差は小さく、高緯度の放出が低緯度より大きい。。低緯度の日射吸収は大きく、正味収支の緯度差を大気・海洋の熱輸送が補う。\par

\noindent\textbf{3：最適} 太陽放射の約半分は可視光線。。太陽光スペクトルの主な一部が可視であるという概数と合う。残りは主に赤外と一部紫外。\par

\noindent\textbf{4：不適切} 入射の約15\%が反射散乱され宇宙へ戻る。。全球反射は約3割の代表値であり、15\%は小さすぎる。\par

\noindent\textbf{5：不適切} 大気下端の太陽垂直面の受光を太陽定数という。。大気の吸収散乱の影響を除いた大気上端で、太陽光に垂直な単位面積あたりの放射束で定義する。\par

\paragraph{再利用する考え方} 全球平均・地域差、貯留量・流量、自然温室効果・追加強制力、代表値・普遍則を区別する。放射の波長帯・反射率・幾何の区別の対象と成立条件を確認する。

\paragraph{共通理論} \hyperref[theory:atmosphere:radiation]{第\ref*{theory:atmosphere:radiation}節}, \hyperref[theory:atmosphere:geometry]{第\ref*{theory:atmosphere:geometry}節}

\par 参考：\cite{atm_nasa_budget,atm_nasa_transport}。

\paragraph{関連問題} \hyperref[q:h26-reg-applied-26]{H26-REG-APPLIED-26}

\par 難易度 2（校準初稿）、優先度未評価。

\subsection{令和元年度再試験 問20}
\label{q:r01-ret-applied-20}

\noindent\textbf{原典：} sources/R01\_retest/17\_applied\_science.pdf、PDF物理7頁。

\textbf{公式正答：} 2（single）。

\paragraph{問題の要旨} 指定された全球平均入射342W/m\textsuperscript{2}に対し、反射散乱で宇宙へ戻る放射束を選ぶ。

\Needspace{4\baselineskip}\paragraph{確認した条件・記述（独自要約）}\mbox{}\par

\noindent\textbf{方向：} 最も適切。\par

\noindent\textbf{指定：} 大気上端の平均342W/m\textsuperscript{2}。宇宙へ戻る反射散乱の和。\par

\paragraph{問われる核心} 全球アルベドと球面平均放射束

\paragraph{必要知識} 全球アルベドと球面平均放射束。問題の採る全球平均値では反射率\(\alpha\simeq0.31\)なので\(F_{\rm refl}=\alpha F_{\rm in}\simeq0.31\times342\simeq106\ {\rm W/m^2}\)。最近の107を選ぶ。342は球面平均であり、ここで再び4で割らない。NASAの代表収支図にも342/107があるが、現在のすべての年の厳密な観測値とはしない。

\paragraph{考え方・正答を選ぶ理由（査読待ち初稿）} 問題の採る全球平均値では反射率\(\alpha\simeq0.31\)なので\(F_{\rm refl}=\alpha F_{\rm in}\simeq0.31\times342\simeq106\ {\rm W/m^2}\)。最近の107を選ぶ。342は球面平均であり、ここで再び4で割らない。NASAの代表収支図にも342/107があるが、現在のすべての年の厳密な観測値とはしない。

\Needspace{4\baselineskip}\paragraph{選択肢の検討（査読待ち初稿）}\mbox{}\par

\noindent\textbf{1：不適切} 68W/m\textsuperscript{2}。入射の約20\%で代表約31\%より小さい。\par

\noindent\textbf{2：最適} 107W/m\textsuperscript{2}。107/342=約0.313で約31\%の全球反射と合う。\par

\noindent\textbf{3：不適切} 137W/m\textsuperscript{2}。約40\%で問題の全球代表反射率より大きい。\par

\noindent\textbf{4：不適切} 171W/m\textsuperscript{2}。約50\%で問題の全球代表反射率より大きい。\par

\noindent\textbf{5：不適切} 205W/m\textsuperscript{2}。約60\%で反射を過大評価している。\par

\paragraph{再利用する考え方} 全球平均・地域差、貯留量・流量、自然温室効果・追加強制力、代表値・普遍則を区別する。全球アルベドと球面平均放射束の対象と成立条件を確認する。

\paragraph{共通理論} \hyperref[theory:atmosphere:radiation]{第\ref*{theory:atmosphere:radiation}節}, \hyperref[theory:atmosphere:geometry]{第\ref*{theory:atmosphere:geometry}節}

\par 参考：\cite{atm_nasa_budget,atm_nasa_albedo}。

\paragraph{関連問題} \hyperref[q:h28-reg-applied-26]{H28-REG-APPLIED-26}

\par 難易度 2（校準初稿）、優先度未評価。

\subsection{令和6年度 問20}
\label{q:r06-reg-applied-20}

\noindent\textbf{原典：} sources/R06/17\_applied\_science.pdf、PDF物理6頁。

\textbf{公式正答：} 3（single）。

\paragraph{問題の要旨} 地球の反射率、太陽放射帯域、季節原因、太陽定数と緯度間の熱収支を判定する。

\Needspace{4\baselineskip}\paragraph{確認した条件・記述（独自要約）}\mbox{}\par

\noindent\textbf{方向：} 最も適切。\par

\paragraph{問われる核心} 短波入射と長波放射の地球規模の収支

\paragraph{必要知識} 短波入射と長波放射の地球規模の収支。反射と吸収、短波日射と長波地球放射を分ける。反射の約31\%が問題採用の適切な概数である。四季は地軸傾斜、太陽定数は大気上端の垂直面、放射不均衡は低緯度から高緯度への輸送につながる。31\%を現在の精密観測で必ず同じとするのではなく、問題の比較精度で扱う。

\paragraph{考え方・正答を選ぶ理由（査読待ち初稿）} 反射と吸収、短波日射と長波地球放射を分ける。反射の約31\%が問題採用の適切な概数である。四季は地軸傾斜、太陽定数は大気上端の垂直面、放射不均衡は低緯度から高緯度への輸送につながる。31\%を現在の精密観測で必ず同じとするのではなく、問題の比較精度で扱う。

\Needspace{4\baselineskip}\paragraph{選択肢の検討（査読待ち初稿）}\mbox{}\par

\noindent\textbf{1：不適切} 地球太陽距離の変化が四季の原因。。地軸の傾きで太陽高度と日長が変わることが主因であり、南北半球の季節は逆になる。\par

\noindent\textbf{2：不適切} 吸収日射の緯度差は小さく、高緯度の放出が低緯度より大きい。。低緯度の日射吸収は大きく、正味収支の緯度差を大気・海洋の熱輸送が補う。\par

\noindent\textbf{3：最適} 入射の約31\%が雲・大気・地表で反射散乱され宇宙へ戻る。。問題で採る全球平均反射率約3割の代表値と合う。時間・雲・地表で変化するので正確な不変定数ではない。\par

\noindent\textbf{4：不適切} 大気下端の太陽垂直面の受光を太陽定数という。。大気の吸収散乱の影響を除いた大気上端で、太陽光に垂直な単位面積あたりの放射束で定義する。\par

\noindent\textbf{5：不適切} 太陽放射の約半分はガンマ線・X線。。太陽放射の主な帯域は可視・近赤外であり、高エネルギー放射が半分ではない。\par

\paragraph{再利用する考え方} 全球平均・地域差、貯留量・流量、自然温室効果・追加強制力、代表値・普遍則を区別する。短波入射と長波放射の地球規模の収支の対象と成立条件を確認する。

\paragraph{共通理論} \hyperref[theory:atmosphere:radiation]{第\ref*{theory:atmosphere:radiation}節}, \hyperref[theory:atmosphere:transport]{第\ref*{theory:atmosphere:transport}節}

\par 参考：\cite{atm_nasa_budget,atm_nasa_albedo,atm_nasa_transport}。

\paragraph{関連問題} \hyperref[q:h26-reg-applied-26]{H26-REG-APPLIED-26}

\par 難易度 2（校準初稿）、優先度未評価。

\section{大気海洋相互作用と南方振動}

\subsection{平成29年度 問21}
\label{q:h29-reg-applied-21}

\noindent\textbf{原典：} sources/H29/17\_applied\_science.pdf、PDF物理7頁。

\textbf{公式正答：} 4（single）。

\paragraph{問題の要旨} 地衡風、境界層、海陸風、ENSO、ロスビー波を時間・空間規模から比較する。

\Needspace{4\baselineskip}\paragraph{確認した条件・記述（独自要約）}\mbox{}\par

\noindent\textbf{方向：} 最も不適切。\par

\paragraph{問われる核心} 循環現象の力学と時間尺度

\paragraph{必要知識} 循環現象の力学と時間尺度。地衡風の力の釣合い、境界層の地表交換、海陸風の日周期、南方振動の年〜数年を分ける。南方振動は太平洋東西の気圧がシーソー状に変化する大気側の現象で、4の数十年は時間尺度が違う。ロスビー波と同じ「振動」という語だけで同じ周期にまとめない。

\paragraph{考え方・正答を選ぶ理由（査読待ち初稿）} 地衡風の力の釣合い、境界層の地表交換、海陸風の日周期、南方振動の年〜数年を分ける。南方振動は太平洋東西の気圧がシーソー状に変化する大気側の現象で、4の数十年は時間尺度が違う。ロスビー波と同じ「振動」という語だけで同じ周期にまとめない。

\Needspace{4\baselineskip}\paragraph{選択肢の検討（査読待ち初稿）}\mbox{}\par

\noindent\textbf{1：適切（条件付き概数）} 地衡風は傾度力とコリオリ力が釣り合い等圧線沿い。。摩擦が小さく加速度が相対的に小さい回転流の近似である。\par

\noindent\textbf{2：適切（条件付き概数）} 地表から約1kmの乱流支配層は大気境界層。。地表の摩擦・熱・水蒸気の交換が強く影響する層の代表厚さ。固定で1kmに鋭い境界があるとはしない。\par

\noindent\textbf{3：適切（条件付き概数）} 海陸風は温度差による日変化の局地風。。海と陸の加熱・冷却の差が水平気圧勾配を作る。\par

\noindent\textbf{4：不適切} 南方振動は熱帯太平洋の気圧の数十年規模東西振動。。ENSOに伴う南方振動は数年規模で、数十年規模の変動とは区別する。\par

\noindent\textbf{5：適切（条件付き概数）} 中緯度偏西風帯の波動はロスビー波。。コリオリ係数の緯度変化が復元性を与える惑星波。中緯度の擾乱と熱輸送に関わる。\par

\paragraph{再利用する考え方} 全球平均・地域差、貯留量・流量、自然温室効果・追加強制力、代表値・普遍則を区別する。循環現象の力学と時間尺度の対象と成立条件を確認する。

\paragraph{共通理論} \hyperref[theory:atmosphere:wind]{第\ref*{theory:atmosphere:wind}節}, \hyperref[theory:atmosphere:local]{第\ref*{theory:atmosphere:local}節}, \hyperref[theory:atmosphere:enso]{第\ref*{theory:atmosphere:enso}節}

\par 参考：\cite{atm_jma_enso,atm_noaa_winds,atm_jma_dynamic}。

\par 難易度 2（校準初稿）、優先度未評価。

\section{降水の酸塩基平衡と酸性沈着}

\subsection{平成27年度 問26}
\label{q:h27-reg-applied-26}

\noindent\textbf{原典：} sources/H27/17\_applied\_science.pdf、PDF物理10頁。

\textbf{公式正答：} 2（single）。

\paragraph{問題の要旨} 降水のpHと水素イオン濃度、炭酸平衡・中和、当時の日本の酸性度を判定する。

\Needspace{4\baselineskip}\paragraph{確認した条件・記述（独自要約）}\mbox{}\par

\noindent\textbf{方向：} 最も不適切。\par

\paragraph{問われる核心} 対数尺度と降水の酸塩基平衡

\paragraph{必要知識} 対数尺度と降水の酸塩基平衡。pHは厳密にはH+活量の負対数であり、希薄溶液ではモル濃度に近似する。\([{\rm H}^+]\simeq10^{-3.3}=5.01\times10^{-4}\ {\rm mol/L}\)なので約501μmol/Lとなり2が不適切。2008〜2012の観測を出題当時の根拠として扱い、現在すべての地点で同じという説明へ流用しない。

\paragraph{考え方・正答を選ぶ理由（査読待ち初稿）} pHは厳密にはH+活量の負対数であり、希薄溶液ではモル濃度に近似する。\([{\rm H}^+]\simeq10^{-3.3}=5.01\times10^{-4}\ {\rm mol/L}\)なので約501μmol/Lとなり2が不適切。2008〜2012の観測を出題当時の根拠として扱い、現在すべての地点で同じという説明へ流用しない。

\Needspace{4\baselineskip}\paragraph{選択肢の検討（査読待ち初稿）}\mbox{}\par

\noindent\textbf{1：適切（条件付き概数）} CO2のみでも25℃の降水pHは約5.6。。CO2と平衡する希薄水の代表的目安であり、全降水の固定値ではない。\par

\noindent\textbf{2：不適切} pH3.3は約30μmol/LのH+。。希薄近似で10\textasciicircum{}-3.3 mol/L=約501μmol/L。30μmol/LならpH約4.52である。\par

\noindent\textbf{3：適切（条件付き概数）} pH5はpH6の10倍のH+。。対数が1違うと濃度比は10となる。pH値を単純に比で比較しない。\par

\noindent\textbf{4：適切（条件付き概数）} 日本には年平均pH5未満の地点が多い。。気象庁が紹介する2008〜2012調査は地点年平均4.48〜5.37、全地点5年4.72で、原典当時の定性的傾向と整合する。全地点5未満ではない。\par

\noindent\textbf{5：適切（条件付き概数）} 乾燥地ではアルカリ粒子・アンモニア中和でpH7に近い場合がある。。土壌の炭酸塩等やNH3は酸を中和する。地域の組成次第で、全乾燥地で固定7とはしない。\par

\paragraph{再利用する考え方} 全球平均・地域差、貯留量・流量、自然温室効果・追加強制力、代表値・普遍則を区別する。対数尺度と降水の酸塩基平衡の対象と成立条件を確認する。

\paragraph{共通理論} \hyperref[theory:atmosphere:acid]{第\ref*{theory:atmosphere:acid}節}

\par 参考：\cite{atm_jma_acid,atm_jma_acid2014}。

\paragraph{関連問題} \hyperref[q:r05-reg-applied-20]{R05-REG-APPLIED-20}

\par 難易度 2（校準初稿）、優先度未評価。

\subsection{令和5年度 問20}
\label{q:r05-reg-applied-20}

\noindent\textbf{原典：} sources/R05/17\_applied\_science.pdf、PDF物理7頁。

\textbf{公式正答：} 1（single）。

\paragraph{問題の要旨} 降水pHの地域差、対数比、炭酸平衡、中和と酸性沈着の影響を判定する。

\Needspace{4\baselineskip}\paragraph{確認した条件・記述（独自要約）}\mbox{}\par

\noindent\textbf{方向：} 最も不適切。\par

\paragraph{問われる核心} 降水の酸塩基平衡と酸性沈着の対象

\paragraph{必要知識} 降水の酸塩基平衡と酸性沈着の対象。海に接している雨と海水は別の相であり、海水のpHをそのまま降水へ代入できない。CO2平衡・硫酸硝酸・塩基性土壌成分等の収支で判断するため1が不適切。pH5.6を人為汚染がない全降水の固定値にもせず、中和と自然酸性源を考える。

\paragraph{考え方・正答を選ぶ理由（査読待ち初稿）} 海に接している雨と海水は別の相であり、海水のpHをそのまま降水へ代入できない。CO2平衡・硫酸硝酸・塩基性土壌成分等の収支で判断するため1が不適切。pH5.6を人為汚染がない全降水の固定値にもせず、中和と自然酸性源を考える。

\Needspace{4\baselineskip}\paragraph{選択肢の検討（査読待ち初稿）}\mbox{}\par

\noindent\textbf{1：不適切} 海洋上・極域の降水はpH8以上が多い。。海水の弱アルカリ性を降水へ移せない。雨のpHはCO2と酸性/中和成分で決まり、場所名だけで一般に8以上としない。\par

\noindent\textbf{2：適切（条件付き概数）} pH5はpH6の10倍のH+。。10\textasciicircum{}-5/10\textasciicircum{}-6=\allowbreak{}10。対数尺度の比較である。\par

\noindent\textbf{3：適切（条件付き概数）} 大気CO2と平衡する水は約pH5.6。。25℃付近の希薄水の代表的目安。CO2分圧や温度、他の溶質が変われば同じ値ではない。\par

\noindent\textbf{4：適切（条件付き概数）} 乾燥地でNH3等中和によりpH7近くなる場合が多い。。アンモニア・土壌由来アルカリ成分は酸を中和する。全乾燥地で必ず7の保証ではない。\par

\noindent\textbf{5：適切（条件付き概数）} 酸性雨は森林・魚類・歴史建造物へ被害を起こす。。酸性沈着による水・土壌の変化と鉱物・金属への作用が生態系と構造物へ影響する。\par

\paragraph{再利用する考え方} 全球平均・地域差、貯留量・流量、自然温室効果・追加強制力、代表値・普遍則を区別する。降水の酸塩基平衡と酸性沈着の対象の対象と成立条件を確認する。

\paragraph{共通理論} \hyperref[theory:atmosphere:acid]{第\ref*{theory:atmosphere:acid}節}

\par 参考：\cite{atm_jma_acid,atm_jma_acid2014}。

\paragraph{関連問題} \hyperref[q:h27-reg-applied-26]{H27-REG-APPLIED-26}

\par 難易度 2（校準初稿）、優先度未評価。

\section{水収支・貯留・平均滞留時間}

\subsection{平成23年度 問20}
\label{q:h23-reg-applied-20}

\noindent\textbf{原典：} sources/H23/17\_applied\_science.pdf、PDF物理7頁。

\textbf{公式正答：} 4（single）。

\paragraph{問題の要旨} 大気中水蒸気の全球平均滞留時間に最も近い日数を選ぶ。

\Needspace{4\baselineskip}\paragraph{確認した条件・記述（独自要約）}\mbox{}\par

\noindent\textbf{方向：} 最も適切。\par

\paragraph{問われる核心} 平均滞留時間を貯留量と流量から定義する

\paragraph{必要知識} 平均滞留時間を貯留量と流量から定義する。定常に近い水の貯留量Wと流出率Fなら平均交換時間の目安は\(\tau=W/F\)である。USGSの大気12900立方km、日蒸発散1170立方kmという別推定の値でも約11日となる。推定法により8〜10日等の差はあるが、本肢では10日が最近。個々の水分子が必ず10日で降るという周期ではない。

\paragraph{考え方・正答を選ぶ理由（査読待ち初稿）} 定常に近い水の貯留量Wと流出率Fなら平均交換時間の目安は\(\tau=W/F\)である。USGSの大気12900立方km、日蒸発散1170立方kmという別推定の値でも約11日となる。推定法により8〜10日等の差はあるが、本肢では10日が最近。個々の水分子が必ず10日で降るという周期ではない。

\Needspace{4\baselineskip}\paragraph{選択肢の検討（査読待ち初稿）}\mbox{}\par

\noindent\textbf{1：不適切} 200日。代表的な約10日に対し約20倍で長すぎる。\par

\noindent\textbf{2：不適切} 120日。約10日の一桁日数とは整合しない。\par

\noindent\textbf{3：不適切} 50日。大気の小さい貯留量と大きい降水流量の比より長すぎる。\par

\noindent\textbf{4：最適} 10日。全球水量を出入りする水量率で割った約8〜11日の代表範囲に最近。\par

\noindent\textbf{5：不適切} 36時間。約1.5日で代表的な全球交換時間より短い。\par

\paragraph{再利用する考え方} 全球平均・地域差、貯留量・流量、自然温室効果・追加強制力、代表値・普遍則を区別する。平均滞留時間を貯留量と流量から定義するの対象と成立条件を確認する。

\paragraph{共通理論} \hyperref[theory:atmosphere:water]{第\ref*{theory:atmosphere:water}節}

\par 参考：\cite{atm_usgs_water,atm_usgs_cycle}。

\par 難易度 2（校準初稿）、優先度未評価。

\subsection{平成28年度 問25}
\label{q:h28-reg-applied-25}

\noindent\textbf{原典：} sources/H28/17\_applied\_science.pdf、PDF物理10頁。

\textbf{公式正答：} 2（single）。

\paragraph{問題の要旨} 水蒸気柱・降水、帯水層係数、氷晶過程、宙水、海陸水収支を比較する。

\Needspace{4\baselineskip}\paragraph{確認した条件・記述（独自要約）}\mbox{}\par

\noindent\textbf{方向：} 最も不適切。\par

\paragraph{問われる核心} 水の貯留と輸送の測定量を分ける

\paragraph{必要知識} 水の貯留と輸送の測定量を分ける。貯留を表す係数と輸送を表す係数を次元で区別する。透水量係数は\(T=Kb\)で単位は面積/時間、比貯留率×厚さは無次元の貯留係数なので2が不適切。可降水量と日雨量も貯留対流量であり、一日の雨が柱の初期貯留を超えても保存則には反しない。

\paragraph{考え方・正答を選ぶ理由（査読待ち初稿）} 貯留を表す係数と輸送を表す係数を次元で区別する。透水量係数は\(T=Kb\)で単位は面積/時間、比貯留率×厚さは無次元の貯留係数なので2が不適切。可降水量と日雨量も貯留対流量であり、一日の雨が柱の初期貯留を超えても保存則には反しない。

\Needspace{4\baselineskip}\paragraph{選択肢の検討（査読待ち初稿）}\mbox{}\par

\noindent\textbf{1：適切（条件付き概数）} 日降水量は可降水量を超えることがある。。可降水量はある時点の柱の貯留で、日降水は時間積分した流量。横からの水蒸気補給があるので上限とはならない。\par

\noindent\textbf{2：不適切} 透水量係数は比貯留率と厚さの積。。T=\allowbreak{}KbでKは透水係数。Ss bは被圧帯水層の貯留係数で、流れやすさとは別量。\par

\noindent\textbf{3：適切（条件付き概数）} 氷晶過程を経る冷たい雨は中緯度降雨に多い。。氷粒子の成長後の融解が雨となる。地上の雨温度が冷たいという命名ではない。\par

\noindent\textbf{4：適切（条件付き概数）} 不飽和帯の低透水層の上に宙水が形成されることがある。。透水性の小さい層が下方への排水を制限し、局所的な地下水面を作り得る。\par

\noindent\textbf{5：適切（条件付き概数）} 海洋は蒸発超過、陸域は降水超過。。長期・全球平均では海洋から陸へ水蒸気が移り、陸から河川等で戻る。個々の地点の常時の関係ではない。\par

\paragraph{再利用する考え方} 全球平均・地域差、貯留量・流量、自然温室効果・追加強制力、代表値・普遍則を区別する。水の貯留と輸送の測定量を分けるの対象と成立条件を確認する。

\paragraph{共通理論} \hyperref[theory:atmosphere:water]{第\ref*{theory:atmosphere:water}節}, \hyperref[theory:atmosphere:rain]{第\ref*{theory:atmosphere:rain}節}, \hyperref[theory:atmosphere:aquifer]{第\ref*{theory:atmosphere:aquifer}節}

\par 参考：\cite{atm_noaa_pwat,atm_mri_cloud,atm_usgs_transmissivity,atm_usgs_perched}。

\par 難易度 2（校準初稿）、優先度未評価。

\subsection{平成29年度 問19}
\label{q:h29-reg-applied-19}

\noindent\textbf{原典：} sources/H29/17\_applied\_science.pdf、PDF物理6頁。

\textbf{公式正答：} 4（single）。

\paragraph{問題の要旨} 潜熱・温室効果、海陸の水収支、氷床の貯留と輸送、平均滞留時間を判定する。

\Needspace{4\baselineskip}\paragraph{確認した条件・記述（独自要約）}\mbox{}\par

\noindent\textbf{方向：} 最も不適切。\par

\paragraph{問われる核心} 水の貯留・交換とエネルギー輸送

\paragraph{必要知識} 水の貯留・交換とエネルギー輸送。交換時間は貯留量/流量である。氷床は大量の水を保持するが長い時間をかけて交換するので、貯留と輸送がともに大きいという肢が不適切。水蒸気の約10日の滞留、潜熱輸送、自然温室効果を同じ分子の別の役割として整理する。水収支の海陸差は地域ごとの恒常的な降水保証ではない。

\paragraph{考え方・正答を選ぶ理由（査読待ち初稿）} 交換時間は貯留量/流量である。氷床は大量の水を保持するが長い時間をかけて交換するので、貯留と輸送がともに大きいという肢が不適切。水蒸気の約10日の滞留、潜熱輸送、自然温室効果を同じ分子の別の役割として整理する。水収支の海陸差は地域ごとの恒常的な降水保証ではない。

\Needspace{4\baselineskip}\paragraph{選択肢の検討（査読待ち初稿）}\mbox{}\par

\noindent\textbf{1：適切（条件付き概数）} 水蒸気の潜熱は低緯度から中緯度の熱輸送に重要。。蒸発時に地表からエネルギーを受け取り移流後の凝結で放出する。\par

\noindent\textbf{2：適切（条件付き概数）} 水蒸気は現在の自然温室効果へ最大に寄与。。自然温室効果への寄与と、人為的な追加放射強制力を区別する。CO2の温暖化上の重要性を否定しない。\par

\noindent\textbf{3：適切（条件付き概数）} 陸は降水超過、海は蒸発超過。。全球長期平均では海から陸への大気輸送と陸から海への流出が釣り合う。\par

\noindent\textbf{4：不適切} 氷床は貯留量も輸送量も大きい。。氷床は大きい貯留だが交換に長い時間を要する。大きい貯留から大きい年間流量を結論できない。\par

\noindent\textbf{5：適切（条件付き概数）} 海洋・水蒸気・地下水・土壌水では水蒸気の滞留が最短。。大気は貯留が小さく出入りが大きいため約10日。地下水等には広い時間分布があり代表平均として比較する。\par

\paragraph{再利用する考え方} 全球平均・地域差、貯留量・流量、自然温室効果・追加強制力、代表値・普遍則を区別する。水の貯留・交換とエネルギー輸送の対象と成立条件を確認する。

\paragraph{共通理論} \hyperref[theory:atmosphere:water]{第\ref*{theory:atmosphere:water}節}, \hyperref[theory:atmosphere:greenhouse]{第\ref*{theory:atmosphere:greenhouse}節}, \hyperref[theory:atmosphere:transport]{第\ref*{theory:atmosphere:transport}節}

\par 参考：\cite{atm_usgs_water,atm_usgs_cycle,atm_jma_ghg2012}。

\paragraph{関連問題} \hyperref[q:r01-reg-applied-21]{R01-REG-APPLIED-21}

\par 難易度 2（校準初稿）、優先度未評価。

\subsection{令和元年度 問21}
\label{q:r01-reg-applied-21}

\noindent\textbf{原典：} sources/R01/17\_applied\_science.pdf、PDF物理7頁。

\textbf{公式正答：} 5（single）。

\paragraph{問題の要旨} 潜熱・温室効果、海陸の水収支、氷床の貯留と輸送、平均滞留時間を判定する。

\Needspace{4\baselineskip}\paragraph{確認した条件・記述（独自要約）}\mbox{}\par

\noindent\textbf{方向：} 最も不適切。\par

\paragraph{問われる核心} 水の貯留・交換とエネルギー輸送

\paragraph{必要知識} 水の貯留・交換とエネルギー輸送。交換時間は貯留量/流量である。氷床は大量の水を保持するが長い時間をかけて交換するので、貯留と輸送がともに大きいという肢が不適切。水蒸気の約10日の滞留、潜熱輸送、自然温室効果を同じ分子の別の役割として整理する。水収支の海陸差は地域ごとの恒常的な降水保証ではない。

\paragraph{考え方・正答を選ぶ理由（査読待ち初稿）} 交換時間は貯留量/流量である。氷床は大量の水を保持するが長い時間をかけて交換するので、貯留と輸送がともに大きいという肢が不適切。水蒸気の約10日の滞留、潜熱輸送、自然温室効果を同じ分子の別の役割として整理する。水収支の海陸差は地域ごとの恒常的な降水保証ではない。

\Needspace{4\baselineskip}\paragraph{選択肢の検討（査読待ち初稿）}\mbox{}\par

\noindent\textbf{1：適切（条件付き概数）} 水蒸気は現在の自然温室効果へ最大に寄与。。自然温室効果への寄与と、人為的な追加放射強制力を区別する。CO2の温暖化上の重要性を否定しない。\par

\noindent\textbf{2：適切（条件付き概数）} 水蒸気の潜熱は低緯度から中緯度の熱輸送に重要。。蒸発時に地表からエネルギーを受け取り移流後の凝結で放出する。\par

\noindent\textbf{3：適切（条件付き概数）} 陸は降水超過、海は蒸発超過。。全球長期平均では海から陸への大気輸送と陸から海への流出が釣り合う。\par

\noindent\textbf{4：適切（条件付き概数）} 海洋・水蒸気・地下水・土壌水では水蒸気の滞留が最短。。大気は貯留が小さく出入りが大きいため約10日。地下水等には広い時間分布があり代表平均として比較する。\par

\noindent\textbf{5：不適切} 氷床は貯留量も輸送量も大きい。。氷床は大きい貯留だが交換に長い時間を要する。大きい貯留から大きい年間流量を結論できない。\par

\paragraph{再利用する考え方} 全球平均・地域差、貯留量・流量、自然温室効果・追加強制力、代表値・普遍則を区別する。水の貯留・交換とエネルギー輸送の対象と成立条件を確認する。

\paragraph{共通理論} \hyperref[theory:atmosphere:water]{第\ref*{theory:atmosphere:water}節}, \hyperref[theory:atmosphere:greenhouse]{第\ref*{theory:atmosphere:greenhouse}節}, \hyperref[theory:atmosphere:transport]{第\ref*{theory:atmosphere:transport}節}

\par 参考：\cite{atm_usgs_water,atm_usgs_cycle,atm_jma_ghg2012}。

\paragraph{関連問題} \hyperref[q:h29-reg-applied-19]{H29-REG-APPLIED-19}

\par 難易度 2（校準初稿）、優先度未評価。

\subsection{令和2年度 問22}
\label{q:r02-reg-applied-22}

\noindent\textbf{原典：} sources/R02/17\_applied\_science.pdf、PDF物理8頁。

\textbf{公式正答：} 1（single）。

\par\textbf{検討未確定：} 公式1を保持。肢4は全球更新と地域枯渇を区別する適用範囲疑義、肢5は日本の流出/蒸発散の比較分母・期間が未指定で定量的な一次統計未確認。

\paragraph{問題の要旨} 蒸発潜熱のエネルギー割合、温室効果、海陸輸送、淡水更新と日本の水収支を比較する。

\Needspace{4\baselineskip}\paragraph{確認した条件・記述（独自要約）}\mbox{}\par

\noindent\textbf{方向：} 最も不適切。\par

\paragraph{問われる核心} 水の再生流量と熱収支の分母

\paragraph{必要知識} 水の再生流量と熱収支の分母。NASA代表収支の潜熱78W/m\textsuperscript{2}と入射342W/m\textsuperscript{2}の比は約23\%である。入射の約1割とする1が不適切で独立解答・公式とも1。ただし4の淡水が枯渇しないという表現は、全球循環の更新と地域資源の不足を区別する必要がある。5も全球又は他国との比較範囲が明示されないので、定量的に全国一律の低蒸発散とは教えない。

\paragraph{考え方・正答を選ぶ理由（査読待ち初稿）} NASA代表収支の潜熱78W/m\textsuperscript{2}と入射342W/m\textsuperscript{2}の比は約23\%である。入射の約1割とする1が不適切で独立解答・公式とも1。ただし4の淡水が枯渇しないという表現は、全球循環の更新と地域資源の不足を区別する必要がある。5も全球又は他国との比較範囲が明示されないので、定量的に全国一律の低蒸発散とは教えない。

\Needspace{4\baselineskip}\paragraph{選択肢の検討（査読待ち初稿）}\mbox{}\par

\noindent\textbf{1：不適切} 入射太陽エネルギーの約1割が蒸発潜熱になる。。代表収支の78/342≈0.23で約1/4、約1割ではない。地表正味日射を分母にすればさらに比が変わる。\par

\noindent\textbf{2：適切} 水蒸気は赤外吸収で温暖な気候へ寄与。。水蒸気は自然温室効果を支える。人為起源の追加強制力と同一のランキングへしない。\par

\noindent\textbf{3：適切（全球長期）} 陸の降水超過は海洋からの水蒸気輸送による。。長期の海陸区分の全球平均収支と整合する。地域や時期の全てで降水超過を保証しない。\par

\noindent\textbf{4：適用範囲に疑義} 淡水資源が枯渇しないのは大気滞留が短く豊富な降水があるため。。全球的に水が再生されることは説明できるが、地域の渇水・地下水過剰取水等の枯渇がないという普遍的保証にはならない。\par

\noindent\textbf{5：定性傾向・定量範囲要確認} 日本は流出が大きく蒸発散が小さい。。日本の多雨と流出優勢という代表的説明。年・流域による変動があり、原典の比較分母や対象期間は未指定である。\par

\paragraph{再利用する考え方} 全球平均・地域差、貯留量・流量、自然温室効果・追加強制力、代表値・普遍則を区別する。水の再生流量と熱収支の分母の対象と成立条件を確認する。

\paragraph{共通理論} \hyperref[theory:atmosphere:water]{第\ref*{theory:atmosphere:water}節}, \hyperref[theory:atmosphere:radiation]{第\ref*{theory:atmosphere:radiation}節}, \hyperref[theory:atmosphere:transport]{第\ref*{theory:atmosphere:transport}節}

\par 参考：\cite{atm_nasa_budget,atm_usgs_cycle}。

\par 難易度 2（校準初稿）、優先度未評価。
